\section{Conclusions}

This study proposes reframing query decomposition and document retrieval for complex query answering as an exploitation–exploration problem under a multi-armed bandit setting. We conduct extensive empirical studies across a range of policies and discover that, when having access to binary relevance labels, Bernoulli distributions with rank information perform best for selecting relevant content under a constrained budget. These results are demonstrated on two datasets, NeuCLIR24 and ResearchyQuestions. We further show that modeling correlations between sub-queries in a hierarchical setting consistently boosts performance across all rewards. We also stress that such an approach performs robustly, while the choice of budget is data-dependent w.r.t.\ the real distribution of relevance across retrieved documents.

We propose further analyses into correlated and contextualized bandits where correlations need not be explicit in the data but may instead emerge naturally in the representational space of sub-queries. Another direction would be to train a retrieval model optimized for relevance, document diversity, exploration, or long-form generation within a policy-learning framework.

\clearpage
\section*{Limitations}
Our work proposes a reframing of an existing paradigm in a popular experimental setup. The study comes with a couple of challenges and limitations: when a retrieval module is involved, the rewards directly depend on the performance of the retriever; moreover, the performance over budgets directly depends on the real distribution of observations. In a dynamic setting, often a set of sub-queries is not pre-defined, such as in the NeuCLIR dataset, which introduces noise by regenerating sub-queries every time. In addition, often the sub-queries are not only not generated, but they are samples from an unbounded search space for a given user request.